\section{Introduction}
The geometry relevant to a prediction problem is the geometry that an observer actually acquires. For a causal experiment this geometry depends on preparation, intervention and stopping, as well as on the physical kernel. A cover of a formally reachable set is not an acquisition procedure. Conversely, a procedure that realizes a favorable conditional distribution may consume many unsuccessful observations before producing its report. The central question is to determine the predictive information that can be acquired and retained by an executable observer, with these costs included.

We take the experiment, not a manifold or an information matrix, as primitive. A preparation and positive report instruments determine future-test evaluations and hence a predictive quotient. A finite observer has a fixed program: it selects an action from its current live label, reads the next report, and either changes live label or enters a terminal label. The latter decision is part of the program. The audit is performed only afterwards. The program cannot read the true model, its own unrecorded past, or an external clock. Persistent labels, program constants, temporary arithmetic, calibration and physical calls are different resources.

Three difficulties arise under an acquisition constraint $\E\tau\le N$. An algebraic balance equation can contain occupation never generated by the preparation. Optimizing the stopping rule changes the acquired law on which compression acts. Finally, approximation of a stopped law in a bounded task metric does not control its mean acquisition cost: a rare discrepancy can initiate an arbitrarily long excursion. A fixed-time quantization bound, or existence of an optimal Borel routing kernel, therefore does not by itself yield an executable finite-description observer under the original budget.

Theorem~\ref{thm:main} connects these three obligations. Its occupation and geometric parts are the causal-stopping results of~\cite{GTFMean}, included with their complete proofs. The constructive part proves that the exact optimum is the limit of certified finite programs using the \emph{same} mean and state budgets, without assuming that all competing observers are uniformly absorbing. This completion starts from uniformly effective unnormalized report instruments, not from an already computed optimum, a smooth routing rule, or a contraction of the normalized posterior. It applies to a known family of experiments with an unknown fixed true model; the program, including its readout, is common to the whole family.

The mechanism is a quantitative conversion of stopping freedom into budget slack. A small fresh probability of routing to an existing terminal label creates a strictly positive mean margin and changes bounded risk by at most that probability times the paid live occupation. No label, clock, or physical reset is added. An occupation-weighted resolvent estimate then separates terminal-risk error from expected-cost error. The former does not acquire an inverse stopping factor, whereas the latter does. This difference permits finite report averaging, probability rounding and certified evaluation to spend only part of the margin. Consequently the finite feasibility test may reject uncertain boundary cases and still contains an approximation to every admissible Borel competitor. The lower certificate is global, not an empirical failure of a search routine.

The same estimates give an explicit same-budget modulus for the optimal value under raw model perturbations. If preparation and one-call instrument errors are $\delta_0$ and $\eta$, then with $\zeta=\delta_0+N\eta$ the sufficient modulus is
\[
 L\{\sqrt{2(N-1)\zeta}+\zeta\}+\alpha,
\]
where $\alpha$ is the terminal-task error; the result is capped by the loss bound $L$. The exponent is not asserted sharp. The theorem includes the exact one-call endpoint $N=1$ separately, rather than dividing by a vanishing budget margin.

For completeness, the inherited occupation argument uses one joint routing inequality over all models. Positivity removes every finite invariant occupation component without changing the terminal ensemble. The mean budget provides compactness, even when other programs never terminate. This proves attainment and finite-subfamily risk obstructions. These statements do not use a model-dependent program, an unrecorded clock, or a formal reachable distribution.

A sharp geometric consequence requires additional information about what a report reveals. Suppose the current task response is recoverable from the current report and the conditional one-call laws are comparable to a reference law $\nu$ that is itself physically prepared. Stopping with mean budget $N$ can increase the density of an acquired terminal law by at most a factor proportional to $N$. The optimal distortion under this density constraint is a mass-constrained quantization profile $\calQ_m^\nu(N)$. A rejection rule realizes the matching upper with one live state and $m$ terminal states. The independent-preparation frontier of the retained adapted-stopping article~\cite{GTFStopped}, reproduced here as a limiting case, gives the exact identities
\begin{equation}\label{eq:intro-exact}
 \calR_{\cp}(N,m)-B=\calQ_m^\nu(N),\qquad
 \calR_{\aut}(N,M)-B=\calQ_{M-1}^\nu(N).
\end{equation}
For correlated controlled reports the profile remains a two-sided comparison. If the acquired reference support is $s$-regular, it yields $(NM)^{-2/s}$ rather than a rate determined by the ambient parameter dimension. There is no requirement that $s$ be an integer. Atoms and rank changes are governed by the profile, not by an artificially assigned common dimension.

The distinction between the two levels is essential. The occupation theorem applies to informative normalized filters with a finite affine closure; the sharp profile theorem does not silently assume that such a filter is report-recoverable. Its classical realization is a genuinely correlated controlled observation experiment, and its quantum realization is a noncommuting instrument with continuously many outcomes. Neither realization is used in the proof of the general theorem. Known-kernel filtering, stationary observation, scheduled horizons and free stopping are separate tasks; a rate in one is not asserted in another without a causal comparison.

\subsection*{Relation to prior work}
Predictive states and minimal future-test representations have an established history; see Shalizi and Crutchfield~\cite{SC} and the experiment-comparison treatment of Fritz, Gonda, Perrone and Rischel~\cite{FGPR}. Finite belief compression is also classical, including the value-directed approach of Poupart and Boutilier~\cite{PB}. These works motivate, but do not supply, the hard-state mean-acquisition theorem proved here.

Occupation equations and constrained control are classical tools. In particular, Dufour and Prieto-Rumeau~\cite{DPR} characterize the invariant part of phantom occupation solutions. We do not claim that decomposition as new. Here it is used for a different purpose: a finite positive predictive representation, a report-only homogeneous controller, and terminal risk under an upper bound on all acquired occupation. One may retain a compact algebraic relaxation and remove its invariant part rather than assume uniform absorption. The proof must still enforce one routing program, including its stopping decision, across all models. Dufour and Prieto-Rumeau's later finite-mixture result~\cite{DPRG} uses uniform absorption for its occupation-space conclusions. Here the finite construction regularizes each competitor within its existing label set, proves a quantitative error against its actual occupation, and pays for the regularization inside the given mean budget. This is not a claim that occupation methods or termination regularization themselves are new. A recent preprint of Sokolov~\cite{Sokolov} gives a Markovization of stopping intensities while preserving observed and surviving measures. A policy of the form $\lambda(t,X_t)$ has access to a state and time; it is not, without further work, a controller with a prescribed finite number of persistent labels.

The functional $\calQ_m^\nu$ is normalized trimmed quantization, not a new static quantization objective. Cuesta-Albertos, Gordaliza and Matr\'an~\cite{CGM} established trimmed $k$-means, and V\'erine, Pydi, Negrevergne and Chevaleyre~\cite{VPNC} study budgeted rejection sampling for distributional objectives. The exact independent-preparation value and its one-live-state realization are inherited from~\cite{GTFStopped}, whose theorem already allows a nonconstant conditional Bayes risk and singular preparation. Those joint-routing and correlated common-program results are inherited from~\cite{GTFMean}. The new step here turns the resulting Borel optimum into arbitrarily accurate finite-description certificates at exactly the same constraints. The small-ball and covering proof is included for every real acquisition budget, so no asymptotic quantization theorem is being invoked outside its hypotheses.

The retained adapted-stopping article~\cite{GTFStopped} also proves a distinct theorem for informative Gaussian and noncommuting weak-effect instruments: at each fixed positive acquisition price it matches the optimal full-history sensing-and-stopping value with a finite-state value. That proof uses calibrated occupation and cost absorption, not report recoverability. We retain it intact rather than identify a Lagrange price with a constrained mean budget without a duality theorem. Its error-composition and calibration mechanisms are inherited tools. Appendix~\ref{app:bayesian} likewise reproduces the earlier Bayesian information decomposition~\cite{GTFBayes}. The retained report-recoverable realizations supply explicit constrained curves. The new informative Gaussian and quantum realizations verify the effective construction on overlapping posteriors, including nonmixing and changing-rank kernels. They provide certified optimal values, not an automatic closed-form power law for these broader classes.

Bayes-risk quantization and Bregman geometry are prior constructions~\cite{Varshney}. Kim's finite-memory study~\cite{Kim} compares fixed-policy costs under one closed-loop execution. Here the converse ranges over all declared adaptive acquisition and stopping policies, while the implementing rule is one common finite-state program. These different quantifiers, and not the mere use of a belief approximation, delimit the contribution.

\subsection*{Scope}
The paper proves a class theorem, not a finite-dimensional reduction of every causal experiment. Finite affine closure, common report domination and the audit interface are explicit in Section~\ref{sec:experiments}. Report recoverability and uniform comparison are additional hypotheses only for the sharp profile. There is no assumption of a preassigned Fisher matrix, no use of total variation to infer a large-deviation principle, and no promotion of finite matrices to unbounded-operator identities. An actual mean budget is not an almost-sure deadline. Under the effective raw certificates the exact constrained optimum has finite two-sided certificates. Efficient synthesis, optimal program length, and a general intrinsic power law for informative filters remain further mathematical questions.
